\documentclass[11pt,a4paper]{article}
\usepackage[T1]{fontenc}
\usepackage{isabelle,isabellesym}

\usepackage{amsmath}
\usepackage{amssymb}
\usepackage{stmaryrd}

% this should be the last package used
\usepackage{pdfsetup}

\urlstyle{rm}
\isabellestyle{it}

\begin{document}

\title{The Relative Arbitrage Problem under Eigenvalue Lower Bounds}
\author{Dmitriy Traytel}
\maketitle

\begin{abstract}
  In a relative-arbitrage problem one asks how quickly a diffusion can be
  driven out of a region when only its instantaneous covariation is
  controlled.  Lai, Shkolnikov and Soner~\cite{LaiShkolnikovSoner} identify
  the answer as the solution of a degenerate elliptic equation; this entry
  formalizes that identification.  For a compact $K \subseteq \mathbb{R}^n$
  and $1 \le k < n$, let $v(x)$ be the supremum, over all martingale laws
  started at $x$ whose instantaneous covariation has its $k$ smallest
  eigenvalues bounded below and its largest bounded above, of the essential
  infimum of the exit time from $K$; we prove all five clauses of their
  Theorem 1.1, that $v$ is finite and upper semicontinuous, that it is a
  viscosity solution of $F(\nabla v, \nabla^{2} v) = 1$ on $K$ with the
  boundary condition of their Definition 3.1, and that it is the only bounded
  upper semicontinuous such solution.  The dynamic programming principle
  their proof defers to Larsson and Ruf~\cite{LarssonRuf} is proved here in
  full, at a deterministic time and at a stopping time, as is the measurable
  selection of an optimizer; the argument runs at a finite horizon, where the
  class of admissible laws is weakly compact, and the paper's own class on
  $C([0,\infty))$ is then shown to have the same value function.  Two points
  of the statement are worth isolating: the boundary condition holds in the
  viscosity sense and not pointwise, since $v > 0$ on the open faces of the
  cube in $\mathbb{R}^{3}$ at $k = 2$, and continuity of $v$ is not asserted
  anywhere.  The reusable machinery --- the second-order differentiability
  theory behind comparison, the construction of Wiener measure, the
  martingale toolbox, and the path-space tightness argument --- is published
  as separate entries.
\end{abstract}

\tableofcontents

\parindent 0pt\parskip 0.5ex

\input{session}

\bibliographystyle{abbrv}
\bibliography{root}

\end{document}
